%% Abstract text, shared by main.tex and sn_main.tex.
Malicious beacon nodes in mobile ad hoc networks corrupt the position
estimates of nodes that never exchange traffic with them, and most existing
defences remove a suspected beacon permanently. We study the decision
before removal and what an error costs. A verification
pipeline fuses behavioural and localisation evidence into per-observer trust
and a three-state verdict; we evaluate it in ns-3 under six isolation
policies and against two literature baselines, in dense, sparse,
shadowed and 50-node networks with up to two attackers. We report four
findings. First, every detector restricted to one evidence class was
blind to the attack outside it, while collecting both classes confirmed the
attacker in 209 of 210 attacked runs of the dense network. Second, we introduce cumulative false
isolation~$\Phi$, in node-seconds: isolating only on confirmation withdrew
3.9--5.0\% of honest-beacon service against 74.8\% for isolating on a
single window, and the gap held when the schedules were enforced. Third, releasing a beacon when its verdict clears re-admitted
forwarding attackers in a closed loop, because isolation silences their
behavioural evidence; releasing only when that evidence returns prevented
every re-admission. Fourth, we derive when weighted Stouffer fusion can
confirm what neither branch confirms alone. The condition explains why, in a
small dense network, a disjunction of the two baselines matched the fused
pipeline, whereas in sparse and larger networks it withdrew 11--24\% of
honest-beacon service against 0--3\% for the pipeline. Evidence coverage and the persistence of errors,
rather than headline accuracy, should guide the evaluation of
beacon verification.
